% ===========================================================================
\chapter{Verificação, Validação e Considerações Finais}
\label{cap_verificacao_validacao}
% ===========================================================================

A integridade e a sustentabilidade de uma arquitetura orientada a objetos dependem da aplicação de métricas formais de qualidade estrutural e de processos contínuos de verificação e validação de software \cite{pressman2021, sommerville2019}. Neste capítulo, a modelagem de classes do SIGEA-GTT é avaliada sob a ótica das métricas clássicas de acoplamento e coesão, da conformidade com os princípios SOLID e do estrito cumprimento do *Quality Gate* institucional de 100\% estabelecido no projeto.

% ===========================================================================
\section{Análise de Métricas Estruturais Orientadas a Objetos}
\label{sec_metricas_oo}
% ===========================================================================

Para quantificar a manutenibilidade do código-fonte e mitigar o surgimento de classes excessivamente complexas (*God Classes*), aplicaram-se as métricas clássicas da suíte de Chidamber e Kemerer:

\begin{enumerate}
    \item \textbf{Acoplamento entre Objetos (\textit{Coupling Between Objects} -- CBO)}:
    \begin{itemize}
        \item Mensura o número de outras classes às quais uma classe está diretamente acoplada.
        \item No SIGEA-GTT, o valor médio de CBO manteve-se baixo ($CBO \le 4$) em razão do uso de Injeção de Dependências e da segregação em camadas. As classes de serviço dependem exclusivamente de interfaces de repositório e mappers, sem criar instâncias diretamente com o operador \texttt{new}.
    \end{itemize}

    \item \textbf{Falta de Coesão nos Métodos (\textit{Lack of Cohesion in Methods} -- LCOM)}:
    \begin{itemize}
        \item Avalia a dispersão de responsabilidades dentro de uma classe medindo o compartilhamento de atributos de instância pelos métodos.
        \item A adesão ao princípio da responsabilidade única (SRP) garantiu alta coesão estrutural ($LCOM \to 0$), demonstrada pelo fato de que cada serviço cuida unicamente de sua respectiva entidade de domínio.
    \end{itemize}

    \item \textbf{Profundidade da Árvore de Herança (\textit{Depth of Inheritance Tree} -- DIT)}:
    \begin{itemize}
        \item Privilegia-se o princípio da composição sobre a herança (\textit{"favor composition over inheritance"}). As entidades JPA não utilizam hierarquias profundas de herança de tabelas, mantendo $DIT \le 2$ e favorecendo a performance e a simplicidade de mapeamento no PostgreSQL.
    \end{itemize}
\end{enumerate}

% ===========================================================================
\section{Cumprimento do Critério de Aceite Estrito (\textit{Quality Gate})}
\label{sec_quality_gate}
% ===========================================================================

Em total conformidade com as diretrizes do projeto departamental e as normas do DCOMP/UFS, a modelagem orientada a objetos assegura o cumprimento dos três pilares inegociáveis do *Quality Gate*:
\begin{enumerate}
    \item \textbf{100\% de Implementação Funcional}: Todas as operações necessárias para dar suporte aos vinte e cinco requisitos funcionais (RF01 a RF25) e aos vinte e dois casos de uso (UC01 a UC22) encontram-se estruturadas e tipadas nas classes e controladores da aplicação;
    \item \textbf{100\% de Cobertura de Testes}: A granularidade das classes e a ausência de acoplamento estático permitem a criação de suítes de testes unitários com JUnit 5 e Mockito no backend, e Jasmine/Karma no frontend, atingindo cobertura de 100\% de linhas e branches auditadas pelo Jacoco;
    \item \textbf{100\% de Documentação Estrutural}: Todos os métodos, atributos, parâmetros e exceções estão aptos à geração automatizada de documentação formal de código por meio do Javadoc e do SpringDoc OpenAPI 3 (Swagger) no backend, e Compodoc no frontend Angular.
\end{enumerate}

% ===========================================================================
\section{Considerações Finais}
\label{sec_consideracoes_finais}
% ===========================================================================

O projeto orientado a objetos e o diagrama de classes consolidado neste documento atestam o alto nível de engenharia e maturidade do sistema SIGEA-GTT. A modelagem em camadas multicamadas (Domínio, Repositório, Serviço, REST e Frontend Reativo com Signals) não apenas atende às exigências acadêmicas do curso de Engenharia de Computação da UFS, mas estabelece uma base sólida, escalável e de fácil manutenção para o avanço da pesquisa multidisciplinar em segurança do paciente e auditoria clínica.
